\documentclass{article}
\usepackage{tcee-repaso}
\tcetema[Comercio y crecimiento]{3.B.9}{Comercio internacional y crecimiento económico. Especial referencia a los efectos del comercio sobre el crecimiento.}
\tcefecha{05/11/2026}
\begin{document}
\cabecerarepaso
\begin{columnas}

% ==========================================================================
\bloque{Esquema del cante}
\esquemacante{%
\apartadocante{Introducción: el huevo o la gallina}{2,5}
\apartadocante{I. Efectos del crecimiento en el comercio}{11}
\subcante{1. Contexto histórico}{1}
\subcante{2. Efectos sobre el comercio}{9}
\subcante{\hspace*{0.8em}País pequeño (\autor{Johnson})}{3}
\subcante{\hspace*{0.8em}País grande (\autor{Bhagwati})}{3,5}
\subcante{\hspace*{0.8em}Acumulación de factores}{1,25}
\subcante{\hspace*{0.8em}Progreso técnico}{1,25}
\subcante{3. Valoración y puente}{1}
\apartadocante{II. Efectos del comercio en el crecimiento}{15,5}
\subcante{Apertura: nivel frente a tasa}{0,75}
\subcante{1. Teoría del crecimiento}{8}
\subcante{\hspace*{0.8em}Los 4 mecanismos}{0,75}
\subcante{\hspace*{0.8em}Exógeno y \autor{Baldwin}}{0,75}
\subcante{\hspace*{0.8em}Externalidades}{1,25}
\subcante{\hspace*{0.8em}\autor{Romer} y \autor{Rivera-Batiz}-\autor{Romer}}{2}
\subcante{\hspace*{0.8em}\autor{Grossman}-\autor{Helpman}}{1,5}
\subcante{\hspace*{0.8em}Otros canales}{0,75}
\subcante{\hspace*{0.8em}Difusión}{1}
\subcante{2. Teoría del desarrollo: ISI e IOE}{2,75}
\subcante{3. Economía geográfica}{1,5}
\subcante{4. Evidencia empírica}{2,5}
\apartadocante{Conclusión}{1}}
\emph{Pizarra} en 3 columnas: izq., $M^A=Y^A_C-Y^A_P$ y tabla de los 5 sesgos (todo el \hyperref[rep:I]{I}); centro, gráficos del \hyperref[rep:I]{I} (5 regiones, ON, empobrecedor, \autor{Rybczynski}), luego $\gamma^{\mathrm E}$ recuadrada; dcha., los 4 mecanismos (todo el \hyperref[rep:II.1]{II.1}). Introducción y conclusión sin pizarra.

% ==========================================================================
\bloque[2,5]{Introducción}
\begin{itemize}
\item \textbf{Enganche}: «¿Qué fue antes, el huevo o la gallina?». Corea del Sur: export. \dato{2,6\,\% del PIB (1960) \ensuremath{\to} 44,4\,\% (2024)}{Banco Mundial, 2026}; PIB por hab. \dato{×36, de 1.038 a 37.048 \$ de 2015}{Banco Mundial, 2026}. ¿Creció porque se abrió o comerció más porque creció?
\item \textbf{Relevancia}: \emph{positiva}, comercio/PIB mundial \dato{26\,\% (1970) \ensuremath{\to} 57\,\% (2024), máx. 62\,\% (2022)}{Banco Mundial, 2026}; 2025, comercio de mercancías \dato{+4,2\,\% frente a PIB +3,0\,\%}{OMC, 2026}. Pequeñas diferencias de tasa $\Rightarrow$ grandes diferencias de renta. \emph{Normativa}: si el comercio cambia la \emph{tasa}, la política comercial es política de crecimiento (aranceles desde 2025, bloques) \vertema{3.B.7}.
\item \textbf{Contextualización}: clásicos ya unían las 2 cosas (\autor{Smith}: mercado y división del trabajo; \autor{Ricardo}: alimentos baratos retrasan el estado estacionario). Luego teoría del comercio \emph{estática} \vertema{3.B.5} \vertema{3.B.6} y del crecimiento \emph{cerrada} \vertema{3.A.43} \vertema{3.A.44}. 2 puentes: años 50 (crecimiento como dato: \autor{Johnson}, \autor{Rybczynski}, \autor{Bhagwati}, \autor{Findlay}-\autor{Grubert}); finales de los 80 (\concepto{ganancias dinámicas}: \autor{Grossman}-\autor{Helpman}, \autor{Rivera-Batiz}-\autor{Romer}, \autor{Young}). En paralelo, \autor{Prebisch}-\autor{Singer} frente a Asia \vertema{3.B.28}.
\item \textbf{Definiciones}: \concepto{crecimiento económico} = aumento sostenido de la renta por hab.; \concepto{grado de apertura} = $(X+M)/Y$.
\item \textbf{Problemática}: ¿cómo afecta el crecimiento, según sesgo y origen, al comercio, la RRI y el bienestar? ¿Puede empobrecer? ¿Cambia el comercio la \emph{tasa}? ¿Qué estrategia y qué dicen los datos?
\item \textbf{Estructura}: \hyperref[rep:I]{I}, crecimiento como dato $\to$ comercio; \hyperref[rep:II]{II} (especial referencia), comercio $\to$ crecimiento. Orden histórico y lógico: el límite de \hyperref[rep:I]{I} (no explica de dónde sale el crecimiento) abre \hyperref[rep:II]{II}.
\end{itemize}

% ==========================================================================
\bloque[11]{I. Efectos del crecimiento en el comercio\label{rep:I}}
\begin{ideaclave}
Sin signo a priori: depende del \emph{sesgo} (consumo y producción), del \emph{origen} (qué factor crece, en qué sector hay progreso) y del \emph{tamaño}. País pequeño: crecer siempre mejora. País grande: el sesgo al exportable empeora la RRI y puede empobrecer.
\end{ideaclave}

\sub{I.1. Contexto histórico\label{rep:I.1}}
\begin{itemize}
\item \autor{Solow} (1956), tras \autor{Harrod} y \autor{Domar} \vertema{3.A.43}: economía \emph{cerrada} y rendimientos decrecientes $\Rightarrow$ estado estacionario, $s\,f(k^*)=(n+\delta)\,k^*$: el comercio no cambia el crecimiento a largo plazo.
\item Por eso se empieza al revés: desde los 50, \emph{estática comparativa} en HO con el crecimiento como dato. \autor{Johnson} (1955, 1958): sesgos; \autor{Rybczynski} (1955): factores; \autor{Bhagwati} (1958): empobrecedor; \autor{Findlay}-\autor{Grubert} (1959): progreso técnico.
\end{itemize}

\sub{I.2. Efectos sobre el comercio\label{rep:I.2}}
\begin{itemize}
\item \emph{Marco}: HO 2×2×2 \vertema{3.B.5}; $A$ abundante en $K$, $X$ intensivo en $K$: $A$ exporta $X$, importa $Y$. Especialización incompleta, preferencias regulares (CIS). Crecer = FPP hacia fuera; se mide por las importaciones a precios constantes:
  \funcionclave{M^A=Y^A_C-Y^A_P}
  Comercio equilibrado ($\mathrm{RRI}\cdot E^A=M^A$): lo que pasa a $M^A$ pasa a $E^A$.
\end{itemize}

\textbf{a) País pequeño (\autor{Johnson}, 1955)}. RRI dada; crecer = $\overline W\uparrow$ (renta en unidades de $Y$).
\begin{itemize}
\item \emph{Consumo}: elasticidad renta de $Y$,
  \begin{formula}
  \[\varepsilon_{Y_C,\overline W}=\frac{\partial Y_C/\partial\overline W}{Y_C/\overline W}=\frac{\mathrm{PMgC}_Y}{Y_C/\overline W}\]
  \end{formula}
  Recta de presupuesto paralela; 5 regiones respecto a $C_0$ (vertical, horizontal y rayo). $=1$ \concepto{neutral} (homotéticas); $>1$ \concepto{procomercio}; $<1$ \concepto{anticomercio}; \concepto{ultraprocomercio} ($\mathrm{PMgC}_Y>1$, $X$ inferior); \concepto{ultraanticomercio} ($Y$ inferior).
\end{itemize}
\graficorepaso[0.5\linewidth]{crecimiento-consumo}{Las 5 regiones de \autor{Johnson} en el consumo}
\begin{itemize}
\item \emph{Producción}: $\varepsilon_{Y_P,\overline W}$, nomenclatura \emph{al revés} (producir más importable sustituye importaciones). $0<\varepsilon<1$ procomercio; $<0$ ultraprocomercio ($Y_P$ cae); $>1$ anticomercio; $\partial Y_P/\partial\overline W>1$ ultraanticomercio ($X_P$ cae).
\item \emph{Efecto total}:
  \begin{formula}
  \[\varepsilon_{M^A,\overline W}=\frac{Y_C}{M^A}\,\varepsilon_{Y_C,\overline W}-\frac{Y_P}{M^A}\,\varepsilon_{Y_P,\overline W}\]
  \end{formula}
  $>1$ procomercio, $=1$ neutral, $<1$ anticomercio. Mismo sesgo en los 2 lados: inmediato; opuestos: deciden las elasticidades ponderadas.
\item \emph{Implicaciones}: el volumen de comercio crece \emph{salvo en el ultraanticomercio}; el \emph{bienestar sube siempre} (solo se amplían las posibilidades de consumo).
\end{itemize}

\textbf{b) País grande (\autor{Bhagwati}, 1958)}. RRI por \concepto{curvas de oferta neta} (\autor{Marshall}, 1879, sobre la demanda recíproca de \autor{Mill}).
\begin{itemize}
\item \emph{RRI}: salvo el ultraanticomercio, ON$_A\to$ ON$_A'$ hacia fuera $\Rightarrow$ exceso de demanda de $Y$ (y, por \autor{Walras}, de oferta de $X$) $\Rightarrow$ $p_Y/p_X\uparrow$: la RRI de $A$ empeora, más cuanto más procomercio y menos elástica la ON del resto del mundo.
\end{itemize}
\graficorepaso[0.5\linewidth]{ofertas-netas}{Crecimiento procomercio de $A$: ON$_A$ se desplaza y la RRI empeora}
\begin{itemize}
\item \emph{Bienestar}: efecto renta (+) y efecto RRI ($-$ si empeora):
  \begin{formula}
  \[\frac{\mathrm{d}V}{\lambda}=\underbrace{\mathrm{d}\overline W_0}_{\text{renta}}+\underbrace{M^A\,\frac{\mathrm{d}\,\mathrm{RRI}}{\mathrm{RRI}}}_{\text{efecto RRI}}\]
  \end{formula}
  Sale de $\overline W=\mathrm{RRI}\cdot X_P+Y_P$, envolvente e identidad de \autor{Roy}. Lo normal: gana el efecto renta. Si gana el RRI: \concepto{crecimiento empobrecedor} (\emph{immiserizing growth}; lo apuntó \autor{Edgeworth}, 1894). «Ir por lana y volver trasquilado».
\end{itemize}
\graficorepaso[0.88\linewidth]{bienestar-grande}{País grande: el bienestar mejora (izq.) o empobrecedor (dcha.)}
\begin{itemize}
\item \emph{Condiciones} (a la vez): (1) crecimiento muy sesgado al exportable (ultraprocomercio en la producción: acumulación del factor del exportable); (2) demanda exterior \emph{inelástica}; (3) país muy dependiente del comercio ($M^A$ grande frente a $\mathrm{d}\overline W_0$); (4) sin política comercial óptima: con el arancel óptimo no se da (\autor{Bhagwati}, 1968) \vertema{3.B.7}.
\item \emph{Implicaciones}: solo empobrece si hay \emph{distorsión}; también en un país pequeño con arancel si crece el sector protegido (\autor{Johnson}, 1967). Argumento de RRI para el arancel óptimo o el impuesto a la exportación (para el país, no para el mundo). Lectura histórica: \autor{Prebisch}-\autor{Singer} (\hyperref[rep:II.2]{II.2}).
\end{itemize}

\textbf{c) Origen: acumulación de factores (\autor{Rybczynski}, 1955)}. País pequeño, especialización incompleta, consumo neutral.
\begin{itemize}
\item \concepto{Teorema de \autor{Rybczynski}}: a precios de los bienes constantes, si crece un factor (el otro fijo), la producción del bien intensivo en él sube \emph{más que proporcionalmente} y la del otro \emph{cae}. $w$, $r$ fijos: consecuencia, no supuesto. \concepto{Efecto magnificación}:
  \begin{formula}
  \[\hat{\overline K}>\hat{\overline L}=0\;\Longrightarrow\;\hat X>\hat{\overline K}>0>\hat Y\]
  \end{formula}
  Intuición: $K/L$ fijo en cada sector; $X$ necesita trabajo que cede $Y$.
\end{itemize}
\graficorepaso[0.82\linewidth]{rybczynski-fpp}{\autor{Rybczynski}: crece $K$ (factor del exportable, izq.) o $L$ (dcha.)}
\begin{itemize}
\item Crece el factor del \emph{exportable}: \emph{ultraprocomercio}, más comercio y más especialización; si es grande, RRI $\downarrow$: vía típica al empobrecedor. Corea, años 60-90: $K/L\uparrow$ (\autor{Young}, 1995) $\Rightarrow$ acero, barcos, automóviles.
\item Crece el factor del \emph{importable}: \emph{ultraanticomercio}, el comercio cae y, si es grande, RRI $\uparrow$.
\item \concepto{Enfermedad holandesa} (\autor{Corden}-\autor{Neary}, 1982; \emph{The Economist}, 1977, Groninga), país \emph{pequeño}, factor específico: el auge del recurso encoge los demás comerciables por (i) \emph{efecto movimiento de recursos} (atrae trabajo y capital) y (ii) \emph{efecto gasto} (precio de los no comerciables $\uparrow$ = apreciación del tipo de cambio real). «Morir de éxito». \emph{Sin} cambio de la RRI. Con aprendizaje, pérdida permanente (\autor{Krugman}, 1987).
\end{itemize}

\textbf{d) Origen: progreso técnico (\autor{Findlay}-\autor{Grubert}, 1959)}.
\begin{itemize}
\item \emph{Neutral} \vertema{3.A.11}: \autor{Hicks}, $A_tF(K,L)$ (RMST fija para $K/L$ dado); \autor{Solow}, $F(A_tK,L)$ (aumentador del capital; participaciones fijas para $L/Y$ dado); \autor{Harrod}, $F(K,A_tL)$ (aumentador del trabajo; para $K/Y$ dado).
\item \emph{Sesgo} (\autor{Hicks}, 1932): \concepto{ahorrador de capital} si, con $K/L$ dado, sube más la PMg del trabajo; \concepto{ahorrador de trabajo}, al revés. Aumentar $\neq$ sesgar: con $\sigma<1$, el aumentador del capital ahorra capital.
\item Supuestos: país pequeño, progreso en un sector, \emph{no reversión de intensidades} (clave para los signos).
\item 2 efectos: \emph{coste} (siempre a favor del sector que progresa; sube la renta del factor en que es intensivo, como \autor{Stolper}-\autor{Samuelson}) y \emph{ahorro de factores} (como \autor{Rybczynski}: refuerza si ahorra el factor del sector, se opone si ahorra el otro).
\item \emph{Resultado}: neutral en $X$ = acumulación de $K$: ultraprocomercio; en $Y$: ultraanticomercio. Si ahorra el factor del \emph{otro} sector: ambiguo.
\item El origen importa menos que el sector: progreso en el exportable $\Rightarrow$ comercio $\uparrow$ y, si es grande, RRI $\downarrow$ (posible empobrecedor).
\end{itemize}

\sub{I.3. Valoración\label{rep:I.3}}
\begin{itemize}
\item \emph{Aporta}: signo ambiguo; lo llamativo, el empobrecedor de un país grande con demanda inelástica y sin arancel óptimo.
\item \emph{Crítica}: estática comparativa, crecimiento como dato; no explica la acumulación ni el progreso, que el comercio puede causar $\Rightarrow$ puente al \hyperref[rep:II]{II}.
\end{itemize}

% ==========================================================================
\bloque[15,5]{II. Efectos del comercio en el crecimiento\label{rep:II}}
\begin{ideaclave}
Neoclásicos: el comercio sube el \emph{nivel} y acelera la transición. Endógenos: puede subir la \emph{tasa}, pero no siempre ni para todos (ideas, especialización, instituciones). Evidencia: efecto positivo en promedio, no automático.
\end{ideaclave}
\begin{itemize}
\item \concepto{Ganancias estáticas} (especialización, escala, variedad): \emph{nivel}, de una vez. \concepto{Ganancias dinámicas}: \emph{nivel} a medio plazo (acumulación) o \emph{tasa} permanente (innovación). Selección de empresas (\autor{Melitz}, 2003): nivel \vertema{3.B.6}.
\end{itemize}

\sub{II.1. Desde la teoría del crecimiento\label{rep:II.1}}
\begin{itemize}
\item 2 interdependencias (\autor{Acemoglu}, 2009): \emph{tecnológicas} (se adopta la tecnología del líder) y \emph{comerciales}. Conocimiento \emph{no rival} $\Rightarrow$ rendimientos crecientes \vertema{3.A.44}.
\item \textbf{4 mecanismos} (\autor{Grossman}-\autor{Helpman}, 2015): (1) difusión internacional del conocimiento (\emph{spillovers}); (2) efecto escala frente a efecto competencia; (3) innovación en equilibrio general (precios de factores y coste de la I+D); (4) difusión de la tecnología (imitar y adoptar).
\end{itemize}

\textbf{a) Crecimiento exógeno}.
\begin{itemize}
\item \emph{Clásicos}: división del trabajo (\autor{Smith}); \emph{vent for surplus} (\autor{Myint}, 1958); difusión de técnicas; importar alimentos elude los rendimientos decrecientes (\autor{Ricardo}).
\item \emph{\autor{Solow}}: nivel, no tasa. \autor{Baldwin} (1992): \concepto{efecto de acumulación a medio plazo} (bono de crecimiento): productividad $\uparrow$ $\Rightarrow$ rendimiento del capital $\uparrow$ $\Rightarrow$ $k^*\uparrow$; bienestar $<$ producción (hay que ahorrar). Mercado único: dinámicos $>$ estáticos (\autor{Baldwin}, 1989) \vertema{3.B.10}.
\end{itemize}

\textbf{b) Externalidades (\autor{Arrow}, 1962; \autor{Romer}, 1986)}. 1.ª generación: progreso = subproducto no buscado de la inversión (aprendizaje por la práctica), que se desborda. Mecanismo 1. «Todo se pega menos la hermosura».
\begin{itemize}
\item Empresas competitivas, ahorro $s$, $n$, $\delta$:
  \begin{formula}
  \[Y_j=A\,K_j^{\alpha}L_j^{1-\alpha}E^{\eta},\qquad E=\frac{X+M}{Y}\,k\]
  \[\gamma_k=s\,A\Bigl(\frac{X+M}{Y}\Bigr)^{\eta}k^{\alpha+\eta-1}-(n+\delta)\]
  \end{formula}
  $\eta=0$: Cobb-Douglas. $\alpha+\eta<1$: estado estacionario, apertura sube el \emph{nivel}; $\alpha+\eta=1$: $AK$, $\gamma=sA[(X+M)/Y]^{\eta}-(n+\delta)$, \emph{abrirse sube la tasa para siempre}; $\alpha+\eta>1$: explosivo.
\end{itemize}
\pareja{externalidades-apertura}{Externalidades con $\alpha+\eta=1$: abrirse sube $\gamma$}{rivera-batiz-romer}{Integrar 2 economías iguales: recta de la tecnología $\uparrow$}
\begin{itemize}
\item Rendimiento social $>$ privado $\Rightarrow$ se invierte de menos: subvencionar la inversión. Capital humano medio: \autor{Lucas} (1988); milagros asiáticos por aprendizaje con mercados exteriores (\autor{Lucas}, 1993).
\item \emph{Valoración}: filo de la navaja $\alpha+\eta=1$; apertura como caja negra.
\end{itemize}

\textbf{c) \autor{Romer} (1990) y \autor{Rivera-Batiz}-\autor{Romer} (1991)}. 2.ª generación: I+D deliberada, nuevas variedades $N_t$ con patente; hogares de \autor{Ramsey} ($\rho$, $\theta$); $L=L_Y+L_I$. Ideas: $\dot N_t=\eta_I L_I N_t$; $\gamma=\eta_I L_I$.
\begin{itemize}
\item Equilibrio de mercado:
  \funcionclave{\gamma^{\mathrm E}=\frac{\alpha\,\eta_I\,L-\rho}{\alpha+\theta}}
  Sube con $L$ (\concepto{efecto escala}) y $\eta_I$; baja con $\rho$. Plano $(\gamma,r)$: hogares (\autor{Euler}) $r=\rho+\theta\gamma$, creciente; tecnología (libre entrada) $r=\alpha\eta_I L-\alpha\gamma$, decreciente.
\item \emph{Integración} de 2 economías iguales (bienes \emph{e ideas}) = $L\to 2L$: $\gamma^{\mathrm E}=(2\alpha\eta_I L-\rho)/(\alpha+\theta)$, \emph{tasa} $\uparrow$ permanente. Efecto \emph{escala} (mercado doble, patente vale el doble) y efecto \emph{redundancia} (no se inventa 2 veces lo mismo): «no hay que reinventar la rueda». Mecanismos 1 y 2.
\item Solo bienes, sin ideas: efecto menor o nulo; limitar flujos de ideas frena tanto como limitar el comercio \vertema{3.B.10}.
\item \emph{Valoración}: efecto escala sin apoyo empírico (\autor{Jones}, 1995) \vertema{3.A.44}; solo países parecidos (sin ventaja comparativa).
\end{itemize}

\textbf{d) \autor{Grossman}-\autor{Helpman}: ventaja comparativa y coste de la I+D}. Mecanismo 3 (1991, caps. 6-7).
\begin{itemize}
\item \emph{Economía pequeña con 2 factores} (1991): bien de \emph{alta tecnología} (intensivo en capital humano $H$, usa las variedades) y \emph{tradicional} (intensivo en trabajo $L$); la I+D usa sobre todo $H$.
\item Cadena \autor{Stolper}-\autor{Samuelson} \vertema{3.B.5}: arancel al \emph{tradicional} $\Rightarrow$ $\hat w_L>\hat p_{\mathrm T}>0>\hat w_H$ $\Rightarrow$ I+D más barata $\Rightarrow$ crecimiento $\uparrow$. Proteger el de \emph{alta tecnología} $\Rightarrow$ $w_H\uparrow$, I+D más cara $\Rightarrow$ crecimiento $\downarrow$.
\item País que \emph{exporta} bienes intensivos en I+D: arancel sobre lo que importa $\Rightarrow$ crece más. País que los \emph{importa} (casi todos los PED): le conviene abrirse.
\item \emph{Crecer más $\neq$ estar mejor}: costes estáticos del arancel \vertema{3.B.7}; mejor subvencionar la I+D.
\item 2 países, un factor (1990): desplazar el gasto hacia el bien del país con ventaja en I+D \emph{reduce} el crecimiento mundial. Norte-Sur; escalera de calidades (1991): el Norte innova, el Sur imita: ciclo del producto de \autor{Vernon} endógeno \vertema{3.B.6}.
\item \emph{Valoración}: \autor{Stolper}-\autor{Samuelson} dinámico; razón de crecimiento, no de bienestar; la protección rara vez sigue la ventaja en I+D.
\end{itemize}

\textbf{e) Otros canales}. Mecanismo 2 en versión menos optimista.
\begin{itemize}
\item \autor{Young} (1991): aprendizaje que se agota y se desborda dentro del país $\Rightarrow$ el atrasado se especializa en bienes sin aprendizaje; crece $\leq$ que en autarquía (aunque gane bienestar estático); el avanzado crece más. Industria naciente dinámica \vertema{3.B.8}.
\item U invertida (\autor{Aghion} et al., 2005): la competencia empuja a las cercanas a la frontera y desanima a las rezagadas.
\item Selección (\autor{Melitz}); Chile (\autor{Pavcnik}, 2002); India, insumos y nuevos productos (\autor{Goldberg} et al., 2010).
\end{itemize}

\textbf{f) Difusión e imitación (\autor{Barro}-\autor{Sala-i-Martin}, 1997)}. Mecanismo 4.
\begin{itemize}
\item Líder $N_1$ innova; seguidor $N_2<N_1$ imita con coste $\nu(N_2/N_1)$ creciente. Imitar $<$ innovar $\Rightarrow$ el seguidor crece más cuanto más lejos y converge (\emph{convergencia condicional}) \vertema{3.A.45}; el comercio acelera la difusión. Corea: de imitar a innovar.
\end{itemize}
\graficorepaso[0.48\linewidth]{difusion}{Difusión: el seguidor crece más cuanto más lejos del líder}
\begin{itemize}
\item Evidencia: convergencia con la liberalización en la CEE (\autor{Ben-David}, 1993); productividad sube con la I+D de los socios (\autor{Coe}-\autor{Helpman}, 1995).
\item \emph{Valoración}: favorece a los atrasados que saben absorber; cerca de la frontera, solo la innovación propia.
\end{itemize}

\sub{II.2. Desde la teoría del desarrollo\label{rep:II.2}}
\begin{itemize}
\item \textbf{CEPAL}: comercio = \emph{obstáculo} para la periferia primaria. \concepto{Tesis de \autor{Prebisch}-\autor{Singer}} (1949-1950): la RRI de los primarios empeora a largo plazo por la \emph{demanda} (ley de \autor{Engel}; protección agrícola del centro) y la \emph{oferta} (en el centro, sindicatos y poder de mercado retienen la productividad; en la periferia, pasa a precios). Lectura histórica del empobrecedor.
\item \textbf{\autor{Thirlwall}} (1979): crecimiento restringido por la balanza de pagos,
  \begin{formula}
  \[\gamma_Y=\frac{\varepsilon_{X,Y^*}}{\varepsilon_{M,Y}}\,\gamma_{Y^*}\]
  \end{formula}
\item \concepto{ISI}: proteger la industria hasta que sea competitiva (industria naciente \vertema{3.B.8}); también racionar divisas (\autor{Irwin}, 2021). América Latina, India, África, España (50-70). \emph{A favor}: los ricos se industrializaron protegidos (\autor{Hamilton}, \autor{List}; \autor{Chang}, 2002). \emph{En contra}: elegir industrias y búsqueda de rentas (\autor{Krueger}, 1974); industria cara, crisis de balanza de pagos; sin difusión. Datos: tendencia débil y con rupturas (\autor{Harvey} et al., 2010).
\item \concepto{IOE}: industrializarse vendiendo al mundo, de lo intensivo en trabajo hacia arriba; 4 tigres desde los 60 \vertema{3.B.28}. \emph{No discrimina contra la exportación}: tipo de cambio efectivo de la exportación $\geq$ el de la importación, neutralidad de incentivos (\autor{Bhagwati}, 1988). Instrumentos: tipo de cambio realista, insumos importados libres para exportar, apoyos condicionados a exportar. Corea: export. \dato{11\,\% del PIB (1970) \ensuremath{\to} 44\,\% (2024)}{Banco Mundial, 2026; BM, 1993}.
\item \emph{Comparación}: ISI, sesgo contra la exportación, aranceles, cuotas, tipo sobrevalorado $\Rightarrow$ industria cara, rentas, crisis; IOE, sesgo neutral, crecimiento rápido con ahorro, educación y estabilidad. NBER (\autor{Bhagwati}, 1978; \autor{Krueger}, 1978) y BM (1993): mejor la orientación exterior. \emph{Matices}: auge de inversión coordinado por el Estado (\autor{Rodrik}, 1995); transpiración más que inspiración (\autor{Young}, 1995; \autor{Krugman}, 1994); la IOE fue también política industrial (\autor{Juhász}-\autor{Lane}-\autor{Rodrik}, 2024). Hoy, ISI como autonomía estratégica \vertema{3.B.8}.
\end{itemize}

\sub{II.3. Desde la economía geográfica\label{rep:II.3}}
\begin{itemize}
\item \concepto{Nueva geografía económica} (\autor{Krugman}, 1991) \vertema{3.B.6}: \emph{dónde} se instala la actividad. Primera naturaleza (Ruhr) y segunda (atracción mutua); fuerzas centrípetas y centrífugas.
\item \emph{Centro-periferia}: 2 regiones iguales; agricultura inmóvil (numerario); manufacturas con rendimientos crecientes, gasto $\mu$, trabajadores móviles ($\lambda_i$); iceberg $\tau$. Aglomeración: vínculo de demanda y de coste de la vida; dispersión: competencia local y demanda agrícola. Con $\phi_\tau\equiv\tau^{1-\sigma}$ por encima del \emph{punto de ruptura} $\phi^{\mathrm r}_\tau$, concentración; no se dispersa hasta bajar del de \emph{sostenibilidad} $\phi^{\mathrm s}_\tau$: \emph{histéresis} (\emph{tomahawk}).
\item \emph{Con crecimiento}: bienes baratos $\Rightarrow$ aglomeración, ideas baratas $\Rightarrow$ dispersión; la aglomeración acelera a las 2 (\autor{Baldwin}-\autor{Forslid}, 2000). Derrames locales: el mundo crece más, rentas divergen (\autor{Martin}-\autor{Ottaviano}, 1999). Despegue del Norte y luego del Sur (\autor{Baldwin}-\autor{Martin}-\autor{Ottaviano}, 2001). Base de la política regional \vertema{3.B.40}.
\end{itemize}

\sub{II.4. Evidencia empírica\label{rep:II.4}}
\begin{itemize}
\item \emph{Dificultad} (el huevo y la gallina): causalidad inversa (el \hyperref[rep:I]{I}), variables omitidas (instituciones, geografía), medición (comercio/PIB es un resultado, no una política).
\item \emph{Entre países}: casi no hay \emph{milagros} ($\geq$3\,\% anual) sin comercio rápido, y los \emph{desastres} rara vez lo tienen (\autor{Panagariya}, 2004). 5 criterios: PED abiertos \dato{4,5\,\% anual por hab. frente a 0,7\,\% los cerrados}{Sachs-Warner, 1995; 1970-1989}. Crítica: \autor{Rodríguez}-\autor{Rodrik} (2000).
\item \emph{Variables instrumentales}: geografía: \dato{+1 pt de comercio/PIB \ensuremath{\Rightarrow} renta +0,5 a 2 pt}{Frankel-Romer, 1999}; \autor{Alcalá}-\autor{Ciccone} (2004), en PPA. Controlando instituciones, se pierde (\autor{Rodrik}-\autor{Subramanian}-\autor{Trebbi}, 2004); Atlántico (\autor{Acemoglu}-\autor{Johnson}-\autor{Robinson}, 2005). \autor{Feyrer} (2019), transporte aéreo: \dato{elasticidad 0,5; 17\,\% de la variación}{1960-1995}; portacontenedores: \dato{0,3}{\autor{Altomonte} et al., 2018, DT}.
\item \emph{Liberalizaciones}: \dato{crecimiento +1,5 pt; inversión +1,5-2 pt; comercio +5 pt del PIB}{\autor{Wacziarg}-\autor{Welch}, 2008}; positivo en promedio, muy distinto por países (\autor{Irwin}, 2025).
\item \emph{Empresas}: importaciones chinas $\Rightarrow$ innovación en Europa, \dato{14\,\% de la mejora técnica}{\autor{Bloom} et al., 2016}; en EE.\,UU., patentes $\downarrow$ (\autor{Autor} et al., 2020): competencia frente a escala.
\item \emph{Valoración}: casi condición necesaria: las \dato{13 economías con $\geq$7\,\% durante $\geq$25 años}{Comisión sobre el Crecimiento, 2008} aprovecharon la economía mundial; no suficiente.
\end{itemize}

% ==========================================================================
\bloque[1]{Conclusión}
\begin{itemize}
\item \textbf{Recapitulación}: relación en los 2 sentidos. Crecimiento $\to$ comercio según sesgo y origen (empobrecedor posible); comercio $\to$ crecimiento: nivel (neoclásico) o tasa (endógeno: ideas, mercado, competencia, coste de la I+D); IOE mejor que ISI; la integración puede concentrar.
\item \textbf{Relevancia}: el comercio cae más que el PIB en las recesiones: \dato{$-12$,2\,\% (2009), $-5$,3\,\% (2020)}{OMC, 2010 y 2021}; con aranceles, saber si afecta al nivel o a la tasa es saber cuánto cuesta cerrarse.
\item \textbf{Extensiones}: \vertema{3.B.6} \vertema{3.B.7} \vertema{3.B.8} \vertema{3.B.10} \vertema{3.B.18} \vertema{3.A.44} \vertema{3.A.45} \vertema{3.B.28}.
\item \textbf{Opinión}: apertura necesaria, no suficiente: estabilidad, instituciones, capital humano (\autor{Rodríguez}-\autor{Rodrik}, 2000).
\item \textbf{Idea final}: ¿el huevo o la gallina? Ninguno: se alimentan el uno al otro; la política decide si el círculo es virtuoso.
\end{itemize}

% ==========================================================================
\bloque{Autores y fechas}
\hito{1776}{\autor{Smith}: mercado, división del trabajo, \emph{vent for surplus}}
\hito{1817}{\autor{Ricardo}: importar alimentos retrasa el estado estacionario}
\hito{1879}{\autor{Marshall}: curvas de oferta neta}
\hito{1949}{\autor{Prebisch} (y \autor{Singer}, 1950): deterioro de la RRI}
\hito{1955}{\autor{Johnson}: sesgos del crecimiento; \autor{Rybczynski}: teorema}
\hito{1956}{\autor{Solow}: crecimiento exógeno}
\hito{1958}{\autor{Bhagwati}: crecimiento empobrecedor}
\hito{1959}{\autor{Findlay}-\autor{Grubert}: progreso técnico y comercio}
\hito{1962}{\autor{Arrow}: aprendizaje por la práctica}
\hito{1966}{\autor{Vernon}: ciclo del producto}
\hito{1967}{\autor{Johnson}: empobrecedor con arancel en país pequeño}
\hito{1968}{\autor{Bhagwati}: sin distorsión no hay empobrecedor}
\hito{1974}{\autor{Krueger}: búsqueda de rentas}
\hito{1979}{\autor{Thirlwall}: restricción de balanza de pagos}
\hito{1982}{\autor{Corden}-\autor{Neary}: enfermedad holandesa}
\hito{1986}{\autor{Romer}: externalidades del capital}
\hito{1988}{\autor{Lucas}: capital humano; \autor{Bhagwati}: IOE}
\hito{1990}{\autor{Romer}: I+D y variedades; \autor{Grossman}-\autor{Helpman}: 2 países}
\hito{1991}{\autor{Rivera-Batiz}-\autor{Romer}; \autor{Grossman}-\autor{Helpman} (economía pequeña, calidades); \autor{Young}; \autor{Krugman}: centro-periferia}
\hito{1992}{\autor{Baldwin}: efecto de acumulación a medio plazo}
\hito{1993}{BM: el milagro de Asia oriental; \autor{Lucas}: milagros}
\hito{1995}{\autor{Sachs}-\autor{Warner}; \autor{Coe}-\autor{Helpman}; \autor{Young}; \autor{Rodrik}; \autor{Jones}}
\hito{1997}{\autor{Barro}-\autor{Sala-i-Martin}: difusión}
\hito{1999}{\autor{Frankel}-\autor{Romer}: geografía como instrumento}
\hito{2000}{\autor{Rodríguez}-\autor{Rodrik}: crítica; \autor{Baldwin}-\autor{Forslid}}
\hito{2005}{\autor{Aghion} et al.: U invertida}
\hito{2008}{\autor{Wacziarg}-\autor{Welch}: liberalizaciones}
\hito{2016}{\autor{Bloom}-\autor{Draca}-\autor{Van Reenen}: China e innovación}
\hito{2019}{\autor{Feyrer}: transporte aéreo}

% ==========================================================================
\bloque{Datos clave}
\begin{itemize}
\item Corea: export. \dato{2,6\,\% (1960), 11\,\% (1970), 44,4\,\% del PIB (2024)}{Banco Mundial, 2026}; PIB por hab. \dato{1.038 \ensuremath{\to} 37.048 \$ de 2015}{Banco Mundial, 2026}.
\item Comercio/PIB mundial: \dato{26\,\% (1970), 62\,\% (2022), 57\,\% (2024)}{Banco Mundial, 2026}.
\item Comercio de mercancías 2025: \dato{+4,2\,\%; PIB +3,0\,\%}{OMC, 2026}; 2009: \dato{$-12$,2\,\%}{OMC, 2010}; 2020: \dato{$-5$,3\,\%}{OMC, 2021}.
\end{itemize}

% ==========================================================================
\bloque{Preguntas del tribunal}
\begin{enumerate}
\item \textbf{¿Puede un país empobrecerse por crecer?} Sí, si es grande, el crecimiento es muy sesgado al exportable, la demanda exterior es inelástica, depende mucho del comercio y no aplica el arancel óptimo (\autor{Bhagwati}, 1958 y 1968).
\item \textbf{¿Enfermedad holandesa = deterioro de la RRI?} No: país pequeño, reasignación interna (movimiento de recursos y gasto), sin cambio de la RRI (\autor{Corden}-\autor{Neary}).
\item \textbf{¿Por qué \autor{Solow} no da efectos de tasa?} Rendimientos decrecientes; solo nivel y acumulación a medio plazo (\autor{Baldwin}).
\item \textbf{¿Siempre sube la tasa con la integración en \autor{Rivera-Batiz}-\autor{Romer}?} Si viajan las ideas y hay efecto escala; con rendimientos decrecientes en las ideas (\autor{Jones}), solo nivel.
\item \textbf{¿Es bueno un arancel para crecer?} Depende del patrón: si abarata el factor de la I+D (país que exporta bienes intensivos en I+D). Aun así, crecer $\neq$ bienestar.
\item \textbf{¿Por qué funcionó la IOE?} No discrimina contra la exportación; aprendizaje, escala y disciplina; con ahorro, educación y estabilidad (y política industrial).
\end{enumerate}

\begin{trampa}
\begin{itemize}
\item Nomenclatura de \autor{Johnson}: en la \emph{producción} va \emph{al revés} que en el consumo (más $Y_P$ = anticomercio).
\item Ultraprocomercio en el consumo: $\mathrm{PMgC}_Y>1$ ($X$ inferior); el comercio cae \emph{solo} con el ultraanticomercio.
\item País pequeño: el bienestar sube \emph{siempre}; el deterioro de la RRI es necesario, no suficiente, para el empobrecedor.
\item \autor{Rybczynski}: precios de los \emph{bienes} constantes (los de los factores, consecuencia); magnificación $\hat X>\hat{\overline K}>0>\hat Y$.
\item Externalidades del capital físico: \autor{Arrow} y \autor{Romer}, no «modelo de \autor{Lucas}» (test de junio de 2024, pregunta 25); \autor{Lucas} (1988) es capital humano.
\item Tasa $\uparrow$ con la apertura \emph{solo} si $\alpha+\eta=1$; si $<1$, nivel.
\item \autor{Grossman}-\autor{Helpman}: el mecanismo \autor{Stolper}-\autor{Samuelson} es el de la \emph{economía pequeña con 2 factores} (1991), no el AER de 1990 (un factor, 2 países). Test de 2022 (p. 33): exporta bienes intensivos en I+D $\Rightarrow$ le conviene un arancel: verdadera. Test de 2014 (p. 31): protección de sectores intensivos en I+D mejora el crecimiento: se marca por eliminación (leída en sentido amplio). Regla: el arancel ayuda si abarata el factor de la I+D.
\item IOE $\neq$ subvencionar la exportación: neutralidad de incentivos (\autor{Bhagwati}, 1988).
\item \emph{Tomahawk}: ruptura y sostenibilidad; histéresis entre ambos. \autor{Baldwin}-\autor{Forslid}: ideas baratas $\Rightarrow$ dispersión.
\end{itemize}
\end{trampa}

\end{columnas}
\end{document}
